\section{Authentication and Authorization}
Authentication can be intentionally simplified for the assessment, but customer ownership and operation-role boundaries must remain explicit. At minimum:
\begin{itemize}
  \item a customer can only access their own booking data;
  \item operation/admin routes require an elevated role;
  \item secrets and environment-specific values are loaded from configuration rather than committed to source control.
\end{itemize}

\section{Input and Business Validation}
The API rejects invalid quantities, categories that do not belong to the selected concert, unpublished concerts, malformed voucher codes, and any client-provided authoritative pricing information.

\section{Request Correlation}
Each request carries or receives a trace/request identifier. Logs for booking, inventory, voucher, and operation transitions include that identifier so a failed booking can be traced across the complete request path.

\section{Suggested Business Metrics}
\begin{itemize}
  \item \code{booking\_attempt\_total};
  \item \code{booking\_success\_total};
  \item \code{booking\_failure\_total\{reason\}};
  \item \code{booking\_duration\_ms};
  \item \code{inventory\_conflict\_total};
  \item \code{voucher\_exhausted\_total};
  \item \code{idempotency\_hit\_total};
  \item \code{reservation\_expired\_total}.
\end{itemize}

\section{Flash-Sale Protection}
Rate limiting is a resilience control, not a correctness mechanism. Even when a rate limiter is bypassed or unavailable, database constraints and atomic updates must still protect inventory and voucher invariants.
